\clearpage\thispagestyle{empty}
\pdfbookmark[0]{Quellen zur Berechnung des Osterfestes (1800)}{osterfest1800-sources}
\begin{center}\textbf{Quellen zur Berechnung des Osterfestes (1800)}\end{center}
\noindent Diese KI-Übertragung des deutschen Originals in heutiges Deutsch folgt Gauß’ Text in den \textit{Werken}, Band~VI (Göttingen, 1874), S.~73--79, einschließlich der allgemeinen Regeln und der am Schluss abgedruckten handschriftlichen Bemerkung. Der Erstdruck erschien in Franz Xaver von Zachs \textit{Monatlicher Correspondenz zur Beförderung der Erd- und Himmels-Kunde}, Band~2, August 1800, Aufsatz~XV, S.~121--130, unter dem Namen Doctor Gaus in Braunschweig. \href{https://www.digitale-sammlungen.de/de/view/bsb10538594?page=127}{Erstdruck}. \href{https://gdz.sub.uni-goettingen.de/id/PPN236007467}{Historischer deutscher Text in den Werken}.
\par\smallskip\noindent Der Erstdruck gibt auf S.~123 für $a$ 9 mögliche Werte an und führt auf S.~123--124 die Werte 0 bis 18 auf. Die Anmerkung in den \textit{Werken} nennt 19. Bei der Zählung der Schaltjahre schließt der Erstdruck 1700 ausdrücklich aus (S.~126). Die \textit{Werke} beziehen das gegebene Jahr ein, wenn es ein Schaltjahr ist. Für den gregorianischen Kalender bis 1699 gibt der Erstdruck $M=12$, $N=2$ an (S.~128). In den \textit{Werken} steht $M=22$, $N=2$ (S.~78).
\par\smallskip\noindent Eine kurze englische Wiedergabe der drei einleitenden Anweisungen und des Ergebnisses findet sich in Eric S. Roberts’ \textit{Programming Abstractions in C++}, Abbildung~2-18, \href{https://tfetimes.com/wp-content/uploads/2015/04/Prog-in-Cpp.pdf}{Frühjahr 2012}, S.~119, und \href{https://web.stanford.edu/dept/cs_edu/resources/textbook/Reader-Beta-2012.pdf}{Herbst 2012}, S.~118--119 (Aufgabe~8 und Abbildung). Zur Kalenderterminologie und den späteren eigenständigen Regeln siehe Reinhold Bien, \href{https://doi.org/10.1007/s00407-004-0078-5}{\textit{Gauß and Beyond: The Making of Easter Algorithms}}, \textit{Archive for History of Exact Sciences} 58 (2004), S.~439--452, besonders S.~439--442.
\par\smallskip\noindent Übertragung und Satz wurden mit KI erstellt.
\clearpage
